\section{An Illustrative Mechanism and Its Limits}\label{sec:model}

The model separates an observed proxy for the information environment from
financial-message contamination. Its role is to organize hypotheses about
conditional associations. The country panel does not identify its structural
parameters or a causal response to information policy.

\subsection{A bounded information-environment proxy}\label{subsec:state}
Let $M_t\in(0,1)$ denote a normalized information-environment index and
$x_t=\lgt(M_t)$. Consider
\begin{equation}\label{eq:state}
 x_{t+1}=c+\rho x_t+\delta s_t+u_{t+1}.
\end{equation}
Here $c$ is an intercept in logit units, $\rho$ is conditional persistence,
and $s_t$ summarizes other countries' index values. An intercept can be
negative; it is not a count or physical inflow of false messages. The
observed index is a rank, so neither $M_t$ nor $x_t$ measures the share of
financial claims that are false. The process offers a parsimonious description
of persistence in an information environment, motivated by narrative economics
\citep{shiller2017narrative,shiller2019narrative}.

\begin{proposition}[Scalar dynamics conditional on external inputs]\label{prop:state}
With $s_t=0$, $u_t=0$ and $0<\rho<1$, the unique steady state is
$\mu=c/(1-\rho)$, and
$x_t=\mu+\rho^t(x_0-\mu)$. Consequently $M_t=\sigma(x_t)\in(0,1)$ and
converges to $\sigma(\mu)$. With independent mean-zero innovations of
variance $\sigma_u^2$, the stationary latent variance is
$\sigma_u^2/(1-\rho^2)$. A deviation in the latent state has half-life
$h_{1/2}=\ln(2)/[-\ln(\rho)]$, exceeding one observation period precisely
when $\rho>1/2$.
\end{proposition}

This half-life concerns a deviation of the aggregate latent index, not the
survival time of a particular narrative. In the stochastic model,
$\E[\sigma(x)]$ generally differs from $\sigma(\E[x])$. The nonlinear
index also need not inherit the latent state's exact half-life.
For jointly evolving countries with a fixed row-stochastic regional matrix
$W$, the corresponding linear system is
$x_{t+1}=c+(\rho I+\delta W)x_t+u_{t+1}$. Its stability requires the
spectral radius of $\rho I+\delta W$ to be below one. For nonnegative
$\rho,\delta$, that radius equals $\rho+\delta$. Scalar stability alone
therefore does not establish stability of an endogenous regional system.

The distinction between latent and observed persistence can be made explicit.
For a deterministic initial deviation $d_0=x_0-\mu$, the observed deviation is
\begin{equation}\label{eq:index-deviation}
 M_t-\sigma(\mu)=\sigma(\mu+\rho^td_0)-\sigma(\mu).
\end{equation}
Because $0<\sigma'(x)=\sigma(x)\{1-\sigma(x)\}\le1/4$, its absolute
value is at most $\rho^t|d_0|/4$. Close to the steady state it is approximately
$\sigma'(\mu)\rho^td_0$. Thus a common latent persistence can coexist with
different movements in the bounded index: the same latent deviation produces
a smaller index change near either endpoint than near its midpoint. Ranking
countries before taking logits further makes this a statement about relative
information environments, rather than absolute quantities of false content.

Persistence also makes long-run summaries sensitive to estimation error.
Writing $h(\rho)=\log(2)/[-\log(\rho)]$ gives
\begin{equation}\label{eq:half-sensitivity}
 h'(\rho)=\frac{\log(2)}{\rho\{\log(\rho)\}^2}>0,
 \qquad h(\rho)\sim\frac{\log(2)}{1-\rho}\quad\text{as }\rho\uparrow1.
\end{equation}
The half-life therefore becomes highly sensitive near unity. This explains
why the baseline persistence $\Rho$ implies $\HalfLife$ years whereas the
approximation $\RhoBC$ implies a much longer horizon. Transforming a point
estimate is useful for interpretation, but it does not remove uncertainty in
the estimated dynamics. An interval reaching unity cannot be mapped into a
finite upper bound on a stationary half-life.

\subsection{Contamination and financing constraints}\label{subsec:pollution}
Let $p\in[0,1]$ be the unobserved probability that a financial signal is
contaminated. A signal about collateral value $V$ is
$s=V+b\eta d+\varepsilon$, where $b\sim\mathrm{Bernoulli}(p)$,
$\eta\in\{-1,1\}$ has equal probabilities, $d>0$, and $\varepsilon$ has
mean zero and variance $\sigma_\varepsilon^2$. Assume the disturbances and
$V$ are mutually independent, conditional on $p$.

\begin{lemma}[Signal contamination]\label{lem:pollution}
$\E[s\mid V,p]=V$ and
$\mathrm{Var}(s\mid V,p)=\sigma_\varepsilon^2+pd^2$.
\end{lemma}

The lemma is about signal noise conditional on the fundamental. It is not a
formula for posterior uncertainty $\mathrm{Var}(V\mid s,p)$, which also
depends on the prior and the signal likelihood. To obtain a transparent
financing mechanism, assume directly that lenders apply a haircut increasing
in noise variance:
\begin{equation}\label{eq:credit}
 L(p)=\lambda\{\bar V-\gamma(\sigma_\varepsilon^2+pd^2)\},
 \qquad \lambda,\gamma\ge0.
\end{equation}
Restrict attention to $L(p)\ge0$; this is a reduced-form credit rule rather
than a derived Bayesian valuation. An entrepreneur with net worth $W_0$ and
desired investment $I^*$ invests $\min\{I^*,W_0+L(p)\}$. The resulting
shortfall can be written
\begin{equation}\label{eq:harm}
 H(p,F)=\max\{ap+bF-\kappa_0,0\},\qquad a,b\ge0,
\end{equation}
where $a=\lambda\gamma d^2$ and $F$ summarizes financing pressure. This
illustration connects noisy information and collateral constraints
\citep{kiyotaki1997credit,bernanke1999financial,bond2012real}.

\begin{proposition}[Increasing differences]\label{prop:amplification}
For $p_1\le p_2$ and $F_1\le F_2$,
$H(p_2,F_1)-H(p_1,F_1)\le H(p_2,F_2)-H(p_1,F_2)$.
If $ap_2+bF-\kappa_0\le0$, the incremental shortfall is zero.
\end{proposition}

Linking this mechanism to the data requires an additional hypothesis that
financial contamination $p$ increases with the broad index $M$, conditional
on other determinants. This mapping is neither observed nor estimated.
Increasing differences survive a monotone mapping $p(M)$, but the magnitude
and a linear interaction specification do not follow from ranks alone.
Moreover, Proposition~\ref{prop:amplification} does not establish small mean
effects or an ordering of growth-quantile slopes without assumptions about
how fragility, constraints and other shocks populate the outcome distribution.
Those patterns are empirical questions, not consequences of the proposition
by themselves.

\subsection{Aggregation and the meaning of a stress interaction}\label{subsec:aggregation}
Heterogeneity in financing slack connects the individual constraint to a
smooth aggregate relationship. Let $K$ be an entrepreneur's slack threshold,
with distribution function $C_K$ and continuous density $f_K$. Hold this
distribution fixed as $p$ and $F$ vary, assume $\E|K|<\infty$, and let $a,b$
be common nonnegative coefficients. Average shortfall is
\begin{equation}\label{eq:aggregate-harm}
 \overline H(p,F)=\E_K[(ap+bF-K)_+],\qquad (z)_+=\max\{z,0\}.
\end{equation}
The fixed-distribution assumption isolates the constraint mechanism;
selection into borrowing or changing firm composition can create additional
effects that are not represented here.

\begin{proposition}[Aggregate constraint exposure]\label{prop:aggregation}
Under these assumptions,
\begin{equation}\label{eq:aggregate-derivatives}
 \overline H_p=aC_K(ap+bF),\qquad
 \overline H_{pF}=ab f_K(ap+bF)\ge0.
\end{equation}
If growth obeys $y=y^0-\chi\overline H(p,F)+\epsilon$ with $\chi\ge0$,
and $p=p(m)$ is differentiable and nondecreasing in an index coordinate $m$,
then, holding $y^0$ and the distribution of $\epsilon$ fixed,
\begin{equation}\label{eq:growth-cross}
 \frac{\partial y}{\partial m}
 =-\chi a p'(m)C_K(ap(m)+bF),\qquad
 \frac{\partial^2 y}{\partial m\,\partial F}
 =-\chi ab p'(m)f_K(ap(m)+bF)\le0.
\end{equation}
\end{proposition}

The level slope depends on the fraction already constrained. The interaction
depends on the density close to the constraint boundary: a rise in fragility
moves additional entrepreneurs into the region where information noise
restricts investment. If almost nobody is constrained, the level association
can be small. If almost everybody is constrained, shortfall can be large but
its interaction with further stress can be small. Strong amplification is
therefore most plausible where an appreciable mass lies near the boundary.

A local expansion around reference values $(m_0,F_0)$ contains
$y_m(m-m_0)$ and $y_{mF}(m-m_0)(F-F_0)$, alongside linear stress and
own-variable curvature terms. The empirical interaction is a tractable
approximation to this surface. Its magnitude combines the prevalence and
density of constraints, the credit sensitivity $a$, growth conversion $\chi$,
the unobserved mapping $p'(m)$ and the units of both regressors. The panel's
linear projection also averages over heterogeneous observations and need not
equal a structural derivative at a particular point. Consequently the
negative interaction $\FragilityTwo$ is consistent with the mechanism's sign
restriction but cannot separately estimate any of these structural objects.

\paragraph{Why amplification need not imply a steeper lower tail.}
Consider the additional location-scale specification
\begin{equation}\label{eq:location-scale}
 y=\ell(m,F)+s(m,F)\epsilon,\qquad s(m,F)>0,
 \qquad Q_\tau(y\mid m,F)=\ell(m,F)+s(m,F)q_\tau,
\end{equation}
where $\epsilon$ has a fixed continuous distribution independent of $(m,F)$
and $q_\tau$ is its $\tau$ quantile. A financing shortfall may enter
$\ell$ through $-\chi\overline H$ without changing $s$. In that case the
index shifts every conditional quantile by the same amount even if
$\ell_{mF}<0$. More generally, the index-slope difference between two
quantiles is $s_m(q_\tau-q_\nu)$. With $\tau<\nu$, a more negative
lower-tail slope follows from $s_m>0$, an extra dispersion assumption that
increasing differences alone does not supply. The bootstrap contrasts test
whether slopes differ in the fitted quantile specification. Their failure to
distinguish the tails from the median can coexist with a negative
mean-growth interaction with banking stress. This location-scale example
clarifies the distinction without imposing it on the estimated regressions.

\subsection{Feedback as a theoretical possibility}\label{subsec:trap}
Suppose a bounded, increasing function $g:\R\to\R$ describes stress feedback
in latent-state units. Consider
\begin{equation}\label{eq:feedback}
 \psi(x)=c+\rho x+\varphi g(x),\qquad \varphi\ge0.
\end{equation}
Writing $g=G\circ\sigma$ is permissible, but the Lipschitz constant below
is that of the complete composite function on the latent scale.

\begin{proposition}[Uniqueness and possible multiplicity]\label{prop:trap}
If $g$ is globally Lipschitz with constant $L_g$ and
$0\le\rho+\varphi L_g<1$, $\psi$ has a unique globally attracting fixed
point. Multiplicity is possible outside that sufficient condition: with
$c=-3/4$, $\rho=1/2$, $\varphi=2$ and
$g(x)=\min\{\max\{x,0\},1\}$, the fixed points are $-3/2$, $1/2$ and
$5/2$, with stable outer points and an unstable middle point.
\end{proposition}

\begin{samepage}
Failure of the contraction inequality does not prove a bifurcation. A smooth
fold additionally requires a fixed point satisfying $\psi'(x)=1$, with
suitable curvature and transversality conditions. The panel identifies
neither $\varphi$ nor $g$, so the example establishes possibility without
locating countries relative to an escalation threshold.

\subsection{A conditional half-life comparison}\label{subsec:burden}
\begin{proposition}[Comparative statics of the scalar latent steady state]
\label{prop:burden}
For $c>0$ and $0<\rho<1$, the elasticities of
$\mu=c/(1-\rho)>0$ are $\varepsilon_c=1$ and
$\varepsilon_\rho=\rho/(1-\rho)$. Equal infinitesimal proportional reductions
in $\rho$ and $c$ reduce $\mu$ by more under a reduction in persistence precisely when
$\rho>1/2$. For $c\ge0$, $\mu$ is convex in $\rho<1$.
\end{proposition}
\end{samepage}

The common derivative of the logistic transformation cancels when comparing
the two marginal reductions in the deterministic index $\sigma(\mu)$ at
the same baseline. This does not imply the same ranking for stochastic
welfare or a multi-channel burden score. For $c=0$ the proportional elasticity
of $\mu$ is undefined. For $c<0$, reducing $\rho$ raises $\mu$ toward zero
and raises $M$ toward $1/2$; a universal pollution-reduction claim is false.
The observation period must also be fixed: an annual coefficient cannot be
used to rank interventions acting on daily narrative lifetimes.

Let $\eta_\rho=-d\log\rho/dB_\rho$ and
$\eta_c=-d\log c/dB_c$ denote proportional parameter changes per unit of
expenditure, evaluated at positive $c,\rho$. Persistence intervention has the
larger marginal effect on the latent steady state only if
\begin{equation}\label{eq:cost}
 \eta_\rho\frac{\rho}{1-\rho}>\eta_c.
\end{equation}
Equal proportional parameter changes are thus an equal-cost comparison only
when $\eta_\rho=\eta_c$. Neither effectiveness schedule is identified by
the country panel. Corrections, friction and authentication can affect more
than one parameter, further limiting a one-to-one policy interpretation.
